\documentclass[11pt]{article}
\usepackage[a4paper,margin=25mm]{geometry}
\usepackage{amsmath,amssymb,hyperref}
\setlength{\emergencystretch}{3em}
\title{A negative sixth-minus-fourth moment in mixed free chaos\\\large Disproof of Conjecture 00000007746}
\author{ziangni-sys}\date{10 October 2026}
\begin{document}\maketitle
\paragraph{Result.} Let $S_1,S_2$ be two freely independent standard free Brownian motions and put
\[
 X=\frac{S_1(1)+S_2(1)}{2\sqrt2}.
\]
Then $X$ is a mixed first-order Wigner integral and
\[
 \varphi(X^2)=\frac14,\qquad \varphi(X^4)=\frac18,\qquad
 \varphi(X^6)=\frac5{64},\qquad
 \varphi(X^6)-\varphi(X^4)=-\frac3{64}<0.
\]
This contradicts the conjecture's universal nonnegative-difference assertion.
\section*{The source's scope}
Both language versions explicitly describe the integrals in question as mixed integrals over several independent free Brownian motions. We use that parenthetical description, with two drivers. We do not posit an additional meaning of ``radial.'' The source neither excludes chaos order $k=1$ nor requires variance one. The kernel is the real two-component function
\[
 f(t)=\frac{1}{2\sqrt2}\,1_{[0,1]}(t)(1,1),\qquad
 X=I_1(f)=\frac1{2\sqrt2}\sum_{i=1}^2\int_0^1 dS_i(t).
\]
It is square integrable, has both components nonzero, and is self-adjoint. Refuting the sign assertion suffices to refute the conjunction; no alternative normality criterion or minimum formula is claimed.
\section*{Actual free Brownian and Fock realization}
Take $H=L^2(\mathbb R_+;\mathbb R^2)$ and its complexification. On full Fock space
$\mathcal F(H)=\mathbb C\Omega\oplus\bigoplus_{n\ge1}H^{\otimes n}$,
left creation is $L(h)\xi=h\otimes\xi$, and its adjoint removes the first tensor factor with its inner-product coefficient. The field $L(h)+L(h)^*$ is bounded and self-adjoint. The vacuum state is
$\varphi(T)=\langle T\Omega,\Omega\rangle$; in particular
$\varphi(T^*T)=\|T\Omega\|^2\ge0$.
The standard full Fock construction gives
\[
 S_i(t)=L(1_{[0,t]}\otimes e_i)+L(1_{[0,t]}\otimes e_i)^*.
\]
These are free Brownian motions with freely independent driver algebras. This construction and the first-chaos identification are described in Kemp--Nourdin--Peccati--Speicher, Sections 1.2--1.3 and 3.1; see the \href{https://mathweb.ucsd.edu/~tkemp/KNPS-AoP-2012.pdf}{published paper}. Using two orthogonal components gives the two independent drivers in the same construction.

For $h_i=1_{[0,1]}\otimes e_i$, the vectors $h_1,h_2$ are orthonormal. The closed subspace generated by their finite tensor words and the vacuum is invariant under all four operators $L(h_i),L(h_i)^*$. It is exactly the two-letter full Fock space, with orthonormal basis indexed by all finite words in $\{1,2\}$. Thus restricting to it preserves every vacuum moment of $X$; there is no truncation.
\section*{Direct operator-moment calculation}
For a coefficient function $v$ on finite words, creation and annihilation have coordinates
\[
 (L_i v)(\emptyset)=0,\quad
 (L_i v)(j w)=1_{i=j}v(w),\quad
 (L_i^*v)(w)=v(iw).
\]
These are the coordinates of the actual bounded Fock operators: creation sends a basis word $w$ to $iw$, and its adjoint removes that prefix. Let
$R=L_1+L_1^*+L_2+L_2^*$ and $a=1/(2\sqrt2)$, so $X=aR$.
Starting from $v_0=1_{\{\emptyset\}}$, recursively define $v_{n+1}=Rv_n$ using the displayed formulas. Each iterate is supported on words of length at most $n$, and therefore is a finite Fock vector. Its empty-word coordinate equals $\varphi(R^n)$.

A return to the empty word consists of creation and annihilation steps that never remove a nonexistent prefix. For $n=2,4,6$, the possible balanced shapes number $1,2,5$, respectively. Each matched creation--annihilation pair has two choices of letter. Hence the return counts, also checked directly by expanding the coordinate recursion, are
\[
 v_2(\emptyset)=2,\qquad v_4(\emptyset)=8,\qquad
 v_6(\emptyset)=40.
\]
Since $a^2=1/8$, multiplying by $a^n$ gives the three moments and the negative difference stated above. This argument computes moments of the actual mixed integral, rather than prescribing a scalar moment sequence.
\section*{Formal verification and normalization}
Lean defines the vacuum and the two creation and annihilation operators on coefficient functions indexed by all finite Boolean words. It defines operator powers, proves scalar-power compatibility, expands the second, fourth and sixth vacuum coordinates, and proves the real identity $a^2=1/8$ with $a=1/(2\sqrt2)$. The final theorem proves both the exact difference $-3/64$ and its strict negativity. The Hilbert-space completion, free Brownian identification, and first-chaos bridge are provided above; the formal coordinate calculation uses their exact operator formulas.

The omission of a variance condition matters: unit-variance semicircular moments have fourth moment $2$ and sixth moment $5$, whereas scaling multiplies them by different powers. Our example has variance $1/4$. It refutes the source as written; it does not refute a separately normalized statement.

Run \texttt{lake build} in \texttt{lean/}. The project pins Lean 4.19.0 and Mathlib commit
\begin{center}\texttt{c44e0c8ee63ca166450922a373c7409c5d26b00b}\end{center}
Printed theorem audits contain only standard axioms; no admitted result, custom axiom, or unsafe decision is used.
\end{document}
